\chapter{Introduction}
\label{chap-introduction}

\begin{otherlanguage*}{english}
Battery-powered environmental sensor networks are typically deployed for
months or years at a time, which makes energy efficiency a first-order
design constraint. Fixed-interval sampling is simple to implement but
wastes energy whenever the monitored quantity changes slowly relative to
the sampling rate \citep{Alippi2008}.

\section{Motivation}
Adaptive sampling schemes adjust the measurement interval based on observed
signal behaviour, sampling more frequently when the signal is changing
quickly and less frequently when it is stable \citep{Jain2004}. Prior work
has largely evaluated such schemes in simulation rather than on real
long-term deployments.

\section{Contributions}
This thesis makes the following contributions:
\begin{itemize}
  \item An adaptive sampling strategy that estimates local signal
    variability from recent readings and adjusts sampling interval
    accordingly.
  \item A field evaluation across three deployed environmental sensor
    networks with different underlying signal dynamics.
  \item An analysis of the energy-fidelity trade-off, showing where
    adaptive sampling offers the largest gains over fixed-interval
    sampling.
\end{itemize}
\end{otherlanguage*}
